% =============================================================================
% 03_1_hauptteil_firmware.tex — Domäne: Firmware
% =============================================================================
% Dieser Abschnitt dokumentiert die Embedded-Software (Firmware) des
% Geotracker-Geräts. Gliederung gemäss Schulvorgabe:
%   1. Analyse und Design
%   2. Softwaredokumentation
%   3. Reflexion
% =============================================================================

\section{Firmware}
\label{sec:firmware}

% Kurze Einführung in die Firmware-Domäne
% Welcher Mikrocontroller / welches Modul wird verwendet?
% Welche Hauptaufgabe erfüllt die Firmware?
<<Kurze Einführung in die verwendete Hardware (z.\,B.\ ESP32, u-blox GPS-Modul)
und die Aufgabe der Firmware innerhalb des Gesamtsystems.>>

% --- 1. Analyse und Design ---------------------------------------------------
\subsection{Analyse und Design}

% Beschreibe die Analysephase: Anforderungen an die Firmware,
% Auswahl der Hardware-Komponenten, Entscheidungen zur Architektur.
% Füge UML-Diagramme (Zustandsdiagramm, Komponentendiagramm) ein,
% die den Ablauf und die Struktur der Firmware zeigen.

\subsubsection{Anforderungsanalyse}
<<Anforderungen an die Firmware auflisten (z.\,B.\ GPS-Daten auslesen,
Koordinaten per MQTT übertragen, Energiesparmodus).>>

\subsubsection{Systemarchitektur und UML-Diagramme}
% Platzhalter für UML-Diagramme (Zustandsautomat, Komponentendiagramm):
% \begin{figure}[htbp]
%   \centering
%   \includegraphics[width=0.85\textwidth]{uml_firmware_zustand}
%   \caption{Zustandsdiagramm der Firmware}
%   \label{fig:uml_firmware}
% \end{figure}
<<UML-Diagramme (Zustands- und Komponentendiagramm) hier einfügen und erläutern.>>

% --- 2. Softwaredokumentation ------------------------------------------------
\subsection{Softwaredokumentation}

\subsubsection{Beschreibung und Konzepte}
% Erkläre die wichtigsten Softwaremodule, Bibliotheken und Konzepte.
% Wie werden GPS-Daten verarbeitet? Wie funktioniert die MQTT-Übertragung?
<<Hauptmodule der Firmware beschreiben (z.\,B.\ GPS-Parser, MQTT-Client,
Schlafmodus-Verwaltung). Wichtige Codeausschnitte als Listings einfügen.>>

% Beispiel für ein Code-Listing:
% \begin{lstlisting}[caption={GPS-Daten auslesen}, label={lst:gps_read}]
%   void readGPS() {
%     // GPS-Daten vom seriellen Port lesen
%   }
% \end{lstlisting}

\subsubsection{Schwierigkeiten und Lösungsansätze}
% Welche technischen Probleme traten auf? Wie wurden sie gelöst?
<<Technische Herausforderungen dokumentieren (z.\,B.\ GPS-Signalverlust in
Gebäuden, Timing-Probleme bei MQTT-Verbindung) und die gewählten
Lösungsstrategien beschreiben.>>

\subsubsection{Testkonzept}
% Wie wurde die Firmware getestet? (Unit-Tests, Integrationstests, Feldtests)
% Welche Testfälle wurden definiert? Wie wurden die Ergebnisse dokumentiert?
<<Testkonzept und Testergebnisse für die Firmware dokumentieren.
Tabelle mit Testfällen und deren Status (bestanden / fehlgeschlagen) einfügen.>>

% --- 3. Reflexion ------------------------------------------------------------
\subsection{Reflexion}
% Bewerte die eigene Arbeit an der Firmware kritisch.
% Was lief gut? Was würde man anders machen?
% Inwiefern wurden die Anforderungen erfüllt?
<<Persönliche und technische Reflexion zur Firmware-Entwicklung.
Stärken und Schwächen des gewählten Ansatzes benennen.>>
